\section{Woran die Theorie scheitern kann}

\textbf{[SETZUNG]} Eine Theorie, die an nichts scheitern kann, sagt nichts. Die
A-Serie führt daher genau vier Größen, an denen sie widerlegbar ist. Sie sind
Verhältnisse oder Strukturaussagen -- nie reine Absolutwerte, denn Absolutwerte
sind Buchhaltung (A080, A130). Die erste, die $\tau$-Masse, ist dabei die
schärfste: eine echte Vorhersage eines messbaren Werts, weil die
Verhältnisstruktur der Leptonen so eng bindet.

Dass es nur vier sind, ist kein Mangel an Mut, sondern eine Folge der
Projektionskette (A090): nur wenige Größen überleben die Projektion aus der
kompakten Struktur in die beobachtete Welt. Was in der Projektion verschwindet,
ist nicht messbar -- nicht aus Ungeschick, sondern strukturell.

\section{Die vier Prognosegrößen}

\textbf{[K]} \textbf{Erstens: die $\tau$-Masse.}
\[
  m_\tau^\text{pred} = 1776{,}968\ \text{MeV},
\]
aus $(m_e, m_\mu)$ und den Strukturzahlen $(\sqrt2, 2/9)$ (A110). Gegen PDG
$1776{,}86\pm0{,}12$ sind das $+0{,}90\,\sigma$. Eine Messung auf $0{,}05$ MeV
entscheidet: landet der Wert bei $\approx 1776{,}97$, ist das Paar
$(\sqrt2, 2/9)$ auf $10^{-5}$ bestätigt; deutlich darunter, ist es widerlegt.
Das ist der schärfste Punkt der Serie und ihr eigentlicher Prognosekern: die
\textbf{Leptonmassen} sagt die Theorie \emph{sehr wohl} voraus. Anders als bei
$\alpha$ oder $G$ steht hier ein Zahlenwert \emph{vor} der Messung fest -- keine
Buchhaltung, keine Ankerwahl, sondern eine echte Vorhersage, weil die
Verhältnisstruktur die dritte Masse aus zweien erzwingt.

\textbf{[K]} \textbf{Zweitens: die $\xi$-Ordnung der CHSH-Abweichung.} Der
Bell-Wert unterschreitet die Tsirelson-Schranke um einen Rest der Ordnung
$\xi$, also $10^{-4}$ bis $10^{-5}$ (A160). Richtung fest (Verringerung),
Größenordnung fest (nicht anpassbar), Skalierung fest (logarithmisch). Heute
nicht auflösbar; ein CHSH-Wert über der Schranke oder deutlich darunter ohne
Hardware-Erklärung widerlegt die Theorie.

\textbf{[K]} \textbf{Drittens: die fraktale Dimension.}
\[
  D_f = 3 - \xi = 2{,}99986667 .
\]
Eine geometrische Struktur mit einer messbar von $3$ verschiedenen Dimension in
genau dieser Größe -- oder ihr Ausbleiben, wo die Theorie sie verlangt --
entscheidet. $D_f$ ist $\xi$ in geometrischer Kleidung; die Abweichung von $3$
\emph{ist} $\xi$.

\textbf{[K]} \textbf{Viertens: die Grundlänge als Ankerpunkt.} Die Definition
\[
  L_0 = \xi\,\ell_P
\]
setzt die Grundlänge der Theorie in Bezug zur Planck-Länge. Unterhalb dieser
Sub-Planck-Skala macht die Theorie keine Aussage (A130). $L_0$ ist die untere
Grenze der Gültigkeit der relativen Beschreibung; eine Struktur, die unterhalb
noch $\xi$-Verhältnisse zeigte, spräche gegen sie.

\textbf{Der Zahlenwert ist Ankersache, kein Test.} $L_0 = \xi\,\ell_P$ ist der
\emph{einzige} Ankerpunkt, an dem die Theorie an eine absolute Skala gebunden
wird -- und wie jeder Anker kann er je nach Blickwinkel anders gesetzt werden
(gegen die Planck-Länge, gegen die Planck-Zeit, gegen $1/\xi$). \emph{Welcher
Anker} gewählt wird, ist Konvention; falsifizierbar ist allein die
$\xi$-Skalierung: dass die Grenze bei $\xi$ mal der Planck-Skala liegt und nicht
bei einer anderen Potenz. Das gilt hier genau wie beim SI-Wert von $\alpha$
(A130) und $G$ (A145).

\textbf{Beim Zahlenwert selbst ist jedoch Sorgfalt geboten -- er ist nicht
beliebig.} Am Standard-Anker (Planck-Länge) ist
\[
  L_0 = \xi\,\ell_P = \tfrac{4}{30000}\cdot 1{,}616\cdot10^{-35}\,\text{m}
       = 2{,}155\cdot10^{-39}\,\text{m}.
\]
Ein gelegentlich auftauchender Wert $5{,}39\cdot10^{-39}$ m ist dagegen
\emph{keine} zulässige Ankervariante, sondern ein Übertragungsfehler: die
Mantisse $5{,}391$ ist die der Planck-\emph{Zeit}
($t_P = 5{,}391\cdot10^{-44}$ s), und der Wert impliziert
$\xi = 3{,}34\cdot10^{-4}$ statt $4/30000 = 1{,}333\cdot10^{-4}$ -- ein Faktor
$2{,}5$ zu groß. Konvention ist die \emph{Wahl} des Ankers; der Zahlenwert bei
gegebenem Anker ist es nicht.

\section{Warum es Verhältnisse sein müssen}

\textbf{[B]} Alle vier sind einheitenfrei oder auf $\xi$ zurückführbar. Das ist
kein Zufall, sondern Prinzip: nur Verhältnisse sind intrinsisch (A080). Ein
Absolutwert -- etwa \enquote{$G = 6{,}674\cdot10^{-11}$} -- ist an eine
Einheitenwahl gebunden und kann die Theorie weder bestätigen noch widerlegen,
weil er von der Skala abhängt, die selbst eine Eingabe ist.

Deshalb steht $\alpha$ nicht auf der Liste, obwohl A130 es ausführlich behandelt:
sein Zahlenwert ist Buchhaltung. Und deshalb steht die $\tau$-Masse darauf,
obwohl sie eine Masse ist: geprüft wird nicht ihr Absolutwert, sondern das
\emph{Verhältnis}, das $(\sqrt2, 2/9)$ zwischen den drei Leptonmassen erzwingt.

\section{Was keine Prüfgröße ist}

\textbf{[SETZUNG]} Ausdrücklich \emph{nicht} zur Falsifikation geführt:

Die Übereinstimmung von $\alpha$ mit $1/137$ (Buchhaltung, A130). Der Zahlenwert
von $G$ (Buchhaltung, A145). Die Leiterverhältnisse auf Prozentniveau (zu grob,
Prüfpunkte statt Prognose, A100). Der Neutrinosektor (nicht prüfbar, A150). Die
These \enquote{SM ist Projektion}, solange die Abbildung fehlt (A190).

Diese Ehrlichkeit gehört zur Falsifizierbarkeit: eine Größe, die man nicht
scharf prüfen kann, als Prüfgröße zu führen, täuscht Testbarkeit vor.

\section{Prüfskript}

\begin{itemize}
  \item Die vier Prognosegrößen mit Wert, Vergleichsgröße und Entscheidungsmaß;
    und der Nachweis, dass jede einheitenfrei oder auf $\xi$ zurückführbar ist:
    \texttt{python/A\_Serie\_Skripte/a210\_prognosegroessen.py}
\end{itemize}

\section{Was offen bleibt}

\textbf{Erstens sind drei der vier heute nicht scharf messbar.} Nur die
$\tau$-Masse ist in Reichweite laufender Experimente. Die anderen drei sind
\emph{im Prinzip} prüfbar, aber noch nicht praktisch. Das ist ehrlich zu nennen.

\textbf{Zweitens ist die Vollständigkeit der Liste nicht bewiesen.} Dass genau
vier Größen die Projektion überleben, ist Bestandsaufnahme (A090), keine
Ableitung. Eine fünfte ist nicht ausgeschlossen.

\section{Ersetzt}

Dieses Dokument sammelt die Prüfgrößen aus A090, A110, A120, A130, A160 und
führt sie als geschlossene Liste. Es nimmt auf: Dok.~292 ($\tau$-Prognose),
Dok.~142 (CHSH), Dok.~101 ($L_0$, $D_f$) sowie die Falsifikations- und
No-Go-Aussagen aus Dok.~249, Dok.~074 (No-Go) und Dok.~193. Eingearbeitet:
R52 (Zahlenwert von $L_0$ am Standard-Anker korrigiert; $5{,}39\cdot10^{-39}$ m
als Übertragungsfehler verworfen).

\section{Verwendung von KI-Werkzeugen}

Dieses Dokument ist unter Verwendung KI-gestützter Werkzeuge entstanden. Die
Fragestellungen, die Setzungen und alle inhaltlichen Entscheidungen stammen vom
Autor; die Werkzeuge formulieren aus, rechnen nach und erstellen Prüfskripte.
Die Verantwortung für Inhalt und Fehler liegt vollständig beim Autor. Der
ausführliche Wortlaut steht in A010.
